\section{Transport of a spin insertion}\label{sec:marked}
Fix a tuned background from Proposition~\ref{flow:tuned}.  A singleton
spin insertion at $y$ is transported by the exact marked map, with the
factor $L^d$ included at every step.  Denote its scaled random list by
$u_j=u_{j,y}$, and set
\begin{equation}\label{mark:charge}
 \alpha_j=\frac1{r_j^d}\sum_i \ell_s(u_{j,i}),\qquad
 A_{j,q}=\E\alpha_j^q.
\end{equation}
The initial insertion is $r_0^d\sigma_y$, so $\alpha_0=1$.
The deterministic representative $w_{j,y}$ from
Proposition~\ref{lin:marked} has normalized spin charge one.  All estimates
below are uniform in $y$, including its position within a grid cell.
Let $\mathcal Q$ be the finite set of even source powers fixed in
Section~\ref{sec:norms}, containing $2,4,400$.  It was chosen before
the background moment exponent $p$ and the blocking ratio $L$.
We prove the moment asymptotics for every $q\in\mathcal Q$.

\subsection{Control of the random marked list}
For a random marked list, write $\norm{u}_{v,q}$ for its weighted
individual-slot $L^q$ norm, with the weights of
Proposition~\ref{norm:source}.  The source-background bound is the mixed
bound there; its exponent $p$ is fixed sufficiently large relative to
$q$.  The following statement supplies both that bound and the small
component perpendicular to spin charge.

\begin{proposition}\label{mark:kernel}
For each $q\in\mathcal Q$ and sufficiently small $D>0$, $A_{j,q}>0$
and there are fixed constants $C_*,B_s,C$ such that, at every scale,
\begin{align}
 \norm{u_j}_{v,q}&\le C_*A_{j,q}^{1/q},\label{mark:lq}\\
 \bigl(\E\norm{u_{j,i}}^q\norm{h_{j,m}}^p\bigr)^{1/p}
 &\le (C_*A_{j,q}^{1/q})^{q/p}B_s s_j
       e^{-vq n_i/p-vn_m}.\label{mark:mixed}
\end{align}
For $j\ge1$ one has the stronger estimate
\begin{equation}\label{mark:kernel-small}
 \norm{u_j-\alpha_jw_{j,y}}_{v,q}
       \le C A_{j,q}^{1/q}s_j.
\end{equation}
The notation $\alpha_jw_{j,y}$ places the representative in its marked
principal slot.  The allowed disorder can be chosen uniformly for
the fixed finite collection $\mathcal Q$.
\end{proposition}

\begin{proof}
Put $a_j=A_{j,q}^{1/q}=\norm{\alpha_j}_{L^q}$ and
$k_j=\norm{u_j-\alpha_jw_{j,y}}_{v,q}$.
The source-line projection is bounded in the weighted norm.  Its
complement contracts under the scaled linear transport, and
Proposition~\ref{norm:source} bounds the nonlinear difference by
$CUs_j$ when $U=C_*a_j$.  Since linear transport preserves normalized
spin charge,
\begin{equation}\label{mark:rough-charge}
 |a_{j+1}-a_j|\le C C_*a_js_j,
 \qquad
 k_{j+1}\le\kappa k_j+C C_*a_js_j,
\end{equation}
where the first inequality is the reverse triangle inequality in
$L^q$.  In particular $a_{j+1}>0$ and
$a_{j+1}/a_j=1+O(s_j)$ once disorder is small.  This proves positivity
before any normalization by the next moment is used.

To close~\eqref{mark:lq}, reserve a fixed part of $C_*a_j$ for the line
and the rest for the kernel.  Choose $C_*$ large enough for the bounded
line norm and the initial list.  The strict contraction in the second
inequality of~\eqref{mark:rough-charge} then reproduces the kernel
allowance relative to $a_{j+1}$, after decreasing $D$.
The mixed output estimate in Proposition~\ref{norm:source} is
\[
 U^{q/p}\bigl(B_s s_j/8+C_{L,B}s_j+o(s_j)\bigr).
\]
Choose $B_s$ after the background constants so that its first two terms
leave strict room below $B_s s_j$, and then decrease $D$ so that the
remaining room pays the changes from $U$ to $C_*a_{j+1}$ and from $s_j$
to $s_{j+1}$.  Initially the source is deterministic, so the mixed
bound follows directly from the background moment bound.  This closes
the simultaneous induction for~\eqref{mark:lq}--\eqref{mark:mixed}.

At zero background the first step gives exactly the representative
$w_{1,y}$.  Its perturbation is $O(a_0s_0)$, hence
$k_1\le Ca_1s_1$.  Divide the kernel recursion in
\eqref{mark:rough-charge} by $a_{j+1}s_{j+1}$.  The ratios
$a_j/a_{j+1}=1+O(s_j)$ and $s_j/s_{j+1}=1+O(G_j)$ make its contraction
coefficient strictly smaller than one.  Iteration proves
\eqref{mark:kernel-small}.
\end{proof}

\subsection{Connected moments and the scalar recurrence}
Fix $q\in\mathcal Q$ and abbreviate $A_j=A_{j,q}$.  Use separate slot indices
for all tensor factors and define
\begin{equation}\label{mark:connected-objects}
 S_j=\E[u_j^{\otimes q}],\qquad
 C_{1,j}=\cum(u_j^{\otimes q},h_j),\qquad
 C_{2,j}=\cum(u_j^{\otimes q},h_j,h_j).
\end{equation}
In each cumulant $u_j^{\otimes q}$ is \emph{one random tensor-valued
argument}; the $q$ marked factors are not separate cumulant groups.
Taking the product of their normalized spin charges gives the scalar
coordinate $A_j$ of $S_j$.  All products and cumulants in
\eqref{mark:connected-objects} are therefore defined entry by entry
before summing their indexed projective norms.

\begin{lemma}\label{mark:connected-comparison}
For $j\ge1$, with the bounded formal profile $Z_{j,q,y}^*$ from
Section~\ref{sec:coefficients},
\begin{align}
 C_{2,j}&=O(A_js_j^3),\label{mark:c2}\\
 C_{1,j}&=A_jg_j Z_{j,q,y}^*
       +O\bigl(A_jG_j(s_j+\rho^j)\bigr),\label{mark:c1}\\
 S_j&=A_jw_{j,y}^{\otimes q}+O(A_js_j).\label{mark:s-profile}
\end{align}
Consequently the marked charge moments obey
\begin{equation}\label{mark:scalar-recursion}
 A_{j+1}=A_j\left(1+c_{j,q}(y)g_j+
                    O\bigl(G_j(s_j+\rho^j)\bigr)\right).
\end{equation}
The recurrence also holds at $j=0$, where its error is allowed to be
$O(G_0)$.
\end{lemma}

\begin{proof}
The mixed estimates give initially
$\norm{C_{k,j}}\le CA_js_j^k$ for $k=1,2$.
Apply the source Taylor rule through two background degrees, counting
the displayed background arguments among those degrees.  For $C_2$
the only retained term is homogeneous linear transport: a nonlinear
term in either background output has at least three background
letters, and a nonlinear term in the transported source tensor also
adds a third letter.  The degree-zero transport of the grouped source
tensor is deterministic and linear on that tensor.  Thus no independent
$S_jK_j$ term is present in this cumulant recursion.  By
Proposition~\ref{lin:marked},
\begin{equation}\label{mark:c2-recur}
 C_{2,j+1}=\mathcal L_{2,j}C_{2,j}+O(A_js_j^3),
 \qquad \norm{\mathcal L_{2,j}}\le\kappa.
\end{equation}
The initial grouped source tensor is deterministic, so $C_{2,0}=0$.
Contraction now proves~\eqref{mark:c2}.  More explicitly, for each of the
finitely many exponents $b$ below,
\begin{equation}\label{mark:varying-convolution}
 \sum_{i<j}\kappa^{j-1-i}A_i s_i^b\le C_b A_j s_j^b.
\end{equation}
Indeed~\eqref{mark:rough-charge} gives
$A_i/A_j\le(1+C\sqrt{\bar g})^{j-i}$, and
$s_i/s_j\le(1+B_0\bar g)^{(j-i)/2}$; their product can be absorbed
in the strict contraction by choosing $D$ small.

The corresponding expansion of $C_1$ has its homogeneous term and
terms with two background letters.  To see their structure without
any assumption on the disorder law, set $U=u^{\otimes q}$ in the
moment-cumulant identity.  For example,
\[
 \E[U\otimes h\otimes h]
 =C_2+S\otimes K+S\otimes m\otimes m
       +C_1\otimes m+m\otimes C_1,
\]
with the indicated permutations of indices.  The source-linear,
background-quadratic term and the source-background term covaried with
one background are obtained from this identity by subtracting their
products of expectations.  Their only possible leading inhomogeneous
contribution is a bounded bilinear operation on $S,K$; all the
$C_2$ and $C_1m$ terms are $O(A_js_j^3)$ at this stage.
Consequently
\begin{equation}\label{mark:c1-recur}
 C_{1,j+1}=\mathcal L_{1,j}C_{1,j}
       +\mathcal H_j(S_j,K_j)+O(A_js_j^3),
 \qquad\norm{\mathcal L_{1,j}}\le\kappa,
\end{equation}
where $\mathcal H_j$ is exactly the sum of these degree-two
inhomogeneous Taylor terms.  Since $C_{1,0}=0$, contraction first gives
$C_{1,j}=O(A_jG_j)$.

In the expansion of $S_{j+1}$, the degree-one terms now use $C_{1,j}$
and $S_jm_j$, and the degree-two terms use $S_jK_j$; all other terms
cost $O(A_js_j^3)$.  Since $m_j,K_j=O(G_j)$, taking product spin charge
improves~\eqref{mark:rough-charge} to
\begin{equation}\label{mark:improved-ratio}
 A_{j+1}/A_j=1+O(G_j).
\end{equation}
To identify its coefficient, expand
$u_j=\alpha_jw_{j,y}+(u_j-\alpha_jw_{j,y})$ in the $q$ tensor factors.
Every term except the all-line term contains a kernel factor.  Hölder,
\eqref{mark:lq}, and~\eqref{mark:kernel-small} bound its expectation by
$CA_js_j$, proving~\eqref{mark:s-profile}.

The formal profile satisfies
\[
 Z_{j+1,q,y}^*=\mathcal L_{1,j}Z_{j,q,y}^*
                   +\mathcal H_j(w_{j,y}^{\otimes q},V_j).
\]
Subtract this equation multiplied by $A_{j+1}g_{j+1}$ from
\eqref{mark:c1-recur}.  The replacements of $S_j$ and $K_j$ by their
line parts cost $O(A_jG_j(s_j+\rho^j))$, by
\eqref{mark:s-profile} and~\eqref{flow:cov-profile}.
The change of scalar factors costs $O(A_jG_j^2)$, by
\eqref{mark:improved-ratio} and $g_{j+1}-g_j=O(G_j^2)$.
Contraction and~\eqref{mark:varying-convolution} prove~\eqref{mark:c1}.
Finally insert the three profile comparisons and
\eqref{flow:mean-profile}, with its terminal term removed, into the
Taylor formula for the spin charge of $S_{j+1}$.  Its leading
coefficient is exactly $c_{j,q}(y)$ by definition.  The accumulated
error is that in~\eqref{mark:scalar-recursion}.  Scale zero has bounded
coefficients and errors of order $A_0G_0$, which the geometric term
allows.
\end{proof}

\subsection{Summation of the moments}
\begin{proposition}\label{mark:moments}
For each $q\in\mathcal Q$ and fixed sufficiently small
$D>0$,
\begin{equation}\label{mark:asymptotic}
 A_{j,q}=j^{\binom q2/8+o(1)}\qquad(j\to\infty),
\end{equation}
uniformly in the marked site $y$.  In addition there is a constant
independent of sufficiently small $D$ such that
\begin{equation}\label{mark:second-uniform}
 A_{j,2}\le C(1+B_0j\bar g)^{1/4}\qquad(j\ge0).
\end{equation}
\end{proposition}

\begin{proof}
The error in~\eqref{mark:scalar-recursion} has summable absolute values,
because
\[
 \sum_{j\ge0}G_j(s_j+\rho^j)<\infty,\qquad
 \sum_{j\ge0}g_j^2<\infty.
\]
All factors are positive and uniformly close to one for small $D$.
Taking logarithms therefore gives
\begin{equation}\label{mark:log-sum}
 \log A_{J,q}=\sum_{j<J}c_{j,q}(y)g_j+O_D(1).
\end{equation}
The error is uniform in $y$.
Let $\mu_q=\binom q2 B_0/8$.  Proposition~\ref{coef:averages} gives
$\sum_{j<J}(c_{j,q}(y)-\mu_q)=o(J)$ uniformly in $y$.
Abel summation then gives
$\sum_{1\le j<J}(c_{j,q}(y)-\mu_q)/j=o(\log J)$:
for each $\eta>0$, the partial sums are at most $\eta j+C_\eta$,
and the summation-by-parts remainder is at most
$\eta\log J+O_\eta(1)$.
The coefficients are bounded, so replacing $1/(B_0j)$ by
$g_j=(1+o(1))/(B_0j)$ costs another $o(\log J)$.
This proves~\eqref{mark:asymptotic} from~\eqref{mark:log-sum}.

For the bound uniform in $D$, write
$c_j=c_{j,2}(y)$ and $\mu=B_0/8$.
Uniform Cesaro convergence gives
$|\sum_{i<k}(c_i-\mu)|\le\eta k+C_\eta$.
Since $G_j$ decreases, summation by parts implies
\begin{equation}\label{mark:uniform-abel}
 \sum_{j<J}c_jG_j
 \le(\mu+\eta)\sum_{j<J}G_j+C_\eta\bar g
 \le\left(\frac18+\frac\eta{B_0}\right)
        \log(1+B_0J\bar g)+C_\eta.
\end{equation}
By~\eqref{flow:uniform-ratio}, replacing $G_j$ by $g_j$ costs at most
$\norm{c}_{\infty}\sup_j|g_j/G_j-1|\sum_{j<J}G_j$.
Choose $\eta$ and then the disorder bound so that the total coefficient
of the logarithm is less than $1/4$.
The sums of the Taylor errors and of $g_j^2$ are uniformly bounded
(in fact they tend to zero with $D$).  Applying the logarithmic
estimate proves~\eqref{mark:second-uniform}.
\end{proof}

\subsection{The first moment}
The first charge moment need not be used as a norm.  Instead, the second
moment estimate supplies an absolute bound for its Taylor remainders.
This avoids presuming either its sign or its comparability to one.

\begin{proposition}\label{mark:first}
Uniformly in $j$ and $y$, for sufficiently small $D>0$,
\begin{equation}\label{mark:first-bound}
 \frac1C\le A_{j,1}=\E\alpha_j\le C.
\end{equation}
More precisely, $\sup_{j,y}|A_{j,1}-1|\to0$ as $D\downarrow0$.
\end{proposition}

\begin{proof}
Set $U_j^*=C(1+B_0j\bar g)^{1/8}$.  Equations~\eqref{mark:lq},
\eqref{mark:mixed}, and~\eqref{mark:second-uniform}, with $q=2$, give
all one-source absolute estimates with $A_j$ in the previous proof
replaced by $U_j^*$.  To retain the full background smallness in this
passage, write $Y=\norm{u_{j,i}}$ and
$H=\prod_m\norm{h_{j,m}}^{a_m}$, with $b=\sum_m a_m\le p$.
Cauchy--Schwarz gives
\[
 \E[YH]\le\bigl(\E[Y^2H]\bigr)^{1/2}
               \bigl(\E[H]\bigr)^{1/2}
 \le C U_j^*s_j^b e^{-vn_i-v\sum_m a_mn_m}.
\]
The first factor uses the mixed $q=2$ measure and the second the
background moment bound.  Unused powers in both measures pay any
large-field indicator in the Taylor remainder as in
Proposition~\ref{norm:source}.  Apply the same connected-cumulant recursions to the
single grouped source $u_j$.  They yield
\[
 C_{2,j}=O(U_j^*s_j^3),\quad C_{1,j}=O(U_j^*G_j),\quad
 S_j=A_{j,1}w_{j,y}+O(U_j^*s_j).
\]
The change of its scalar coordinate is bounded absolutely by
$|A_{j+1,1}-A_{j,1}|\le CU_j^*G_j$.
When comparing the $C_1$ recursion with its formal profile, use this
absolute difference bound in place of~\eqref{mark:improved-ratio}.
Together with $|A_{j,1}|\le U_j^*$ it bounds the change of
$A_{j,1}g_j$ by $CU_j^*G_j^2$.  The same contracting convolution proves
\[
 C_{1,j}=A_{j,1}g_jZ_{j,1,y}^*
               +O\bigl(U_j^*G_j(s_j+\rho^j)\bigr).
\]
Consequently, with $a_j=A_{j,1}$,
\begin{equation}\label{mark:first-recursion}
 a_{j+1}=(1+c_{j,1}(y)g_j)a_j+e_j,
 \qquad |e_j|\le CU_j^*G_j(s_j+\rho^j).
\end{equation}

Corollary~\ref{coef:bounded-first} gives a fixed bound on every
partial sum of $c_{j,1}(y)$, uniformly in $y$.  Abel summation and
Corollary~\ref{flow:variance} therefore imply
\[
 \sup_J\left|\sum_{j<J}c_{j,1}(y)g_j\right|
 \le C\left(g_0+\sum_{j\ge0}|g_{j+1}-g_j|\right)
 \le C\bar g.
\]
Also $\sum_jc_{j,1}(y)^2g_j^2\le C\bar g$.
Thus the products
$P_J=\prod_{j<J}(1+c_{j,1}(y)g_j)$ are positive and satisfy
$\sup_{J,y}|P_J-1|\le C\bar g$.
The inhomogeneous errors have total cost
\begin{align*}
 \sum_{j\ge0}|e_j|
 &\le C\bar g^{3/2}\sum_{j\ge0}
             (1+B_0j\bar g)^{-11/8}
       +C\bar g\sum_{j\ge0}\rho^j
             (1+B_0j\bar g)^{-7/8}\\
 &\le C\sqrt{\bar g}.
\end{align*}
The exponent $11/8>1$ is what makes the first sum integrable on the
scale $j\asymp\bar g^{-1}$.
Variation of constants in~\eqref{mark:first-recursion}, with $a_0=1$,
gives
\[
 a_J=P_J\left(1+\sum_{j<J}e_j/P_{j+1}\right).
\]
It follows that $\sup_{J,y}|a_J-1|\le C\sqrt{\bar g}$, which proves
both assertions.
\end{proof}
